\section{One row certifies signed Pell pairs and their indices}
\label{sec:one-row-pell}

The useful change is to make the same nonnegative expression serve as a
size bound and as a congruence test.  A separate upper-bound equation is
unnecessary.  Moreover, a very small residue distinguishes the two signs
of a Pell abscissa, so the same expression can orient both coordinates.
The result below is an ordinary arithmetic theorem.  Its application to
the ZFC equation retains the complete inherited indexed-Pell certificate
and the subsequent arithmetic-source calibration.

For $A\ge2$, write
\[
 (A+\sqrt{A^2-1})^j
   =\chi_A(j)+\psi_A(j)\sqrt{A^2-1}\qquad(j\ge0).
\]
We use the standard integer-coordinate Pell classification and addition
formulas.  In particular, every integer solution of
$X^2-(A^2-1)Y^2=1$ has
\[
 X=\sigma\chi_A(j),\qquad Y=\varepsilon\psi_A(j),
 \quad j\ge0,\quad \sigma,\varepsilon\in\{1,-1\}.
\]
The sign of $Y$ is immaterial when $j=0$.  Reduction of the defining
recurrence modulo $A-1$ gives
\begin{equation}\label{eq:one-row-residues}
 \chi_A(j)\equiv1,\qquad \psi_A(j)\equiv j\pmod{A-1}.
\end{equation}
Also $\psi_A(j)\ge j$, and both sequences are strictly increasing after
their zero endpoints.

\subsection{The exact inherited anchor contract}

One complete indexed-Pell certificate is retained at an index $W$.
Denote its principal pair by $(U,V)$ and its private witnesses by
$\boldsymbol\xi$.  The contract used here is precisely this: under the
certificate's stated domain conditions and $V\ge W\ge1$, its equations
imply
\[
 U=\chi_A(W),\qquad V=\psi_A(W).
\]
Conversely, this true pair has private witnesses satisfying the same
equations and domain conditions.  We remove only a separate witness or
row for $V\ge W$; the new aggregate row proves that inequality before
the anchor theorem is invoked.

For example, a usual nonnegative-witness version of this interface has
the equations
\begin{align*}
 U^2-(A^2-1)V^2&=1,&
 U_1^2-(A^2-1)V_1^2&=1,\\
 S_1^2-(B_0^2-1)R_1^2&=1,&
 B_0&=1+4V(q+1)=A+U_1b,\\
 V_1&=V^2(v+1),&
 S_1&=U+U_1s,\\
 R_1&=W+4Vr.
\end{align*}
All the nonnegativity and positive-multiple conditions of the chosen
literal version remain in force.  The following theorem permits the
\emph{principal seed pairs} to be signed.  It does not silently permit
signed private quotients in this anchor.  If an independently proved
signed anchor is used instead, its complete hypotheses must likewise
be retained.

\begin{lemma}[A square separates the next index]\label{lem:next-index-square}
Let $A\ge2$, $B\ge2$, and $1\le e<B$.  Then
\[
 \bigl(\psi_A(B+1)-e\bigr)^2>\psi_A(2B).
\]
\end{lemma}
\begin{proof}
Put $y=\psi_A(B)$ and $x=\chi_A(B)$.  The norm identity gives
$x\le Ay$, because $y\ge1$, and also $x>(A-1)y$.  Since $y\ge B$,
\[
 \psi_A(B+1)-e=Ay+x-e>(2A-2)y+1.
\]
The addition formula gives $\psi_A(2B)=2xy\le2Ay^2$.  Finally,
$(2A-2)^2\ge2A$ for $A\ge2$; equality at $A=2$ is harmless because
of the additional positive terms after squaring.  Thus
$((2A-2)y+1)^2>2Ay^2$, proving the claim.
\end{proof}

\begin{theorem}[A single row for oriented indexed Pell pairs]
\label{thm:one-row-oriented}
Let $W=2B\ge4$ be even, let $A-1>W^3$, and let
\[
 1\le e_i<B\quad(1\le i\le m),\qquad m\le W-2.
\]
Retain the complete anchor contract above.  For each $i$, impose just
the principal norm
\begin{equation}\label{eq:one-row-principal}
 X_i^2-(A^2-1)Y_i^2=1,
 \qquad X_i,Y_i\in\Z.
\end{equation}
Add the single aggregate equation
\begin{equation}\label{eq:one-row-oriented}
 \boxed{\quad
 V=W+\sum_{i=1}^m\bigl((X_i-1)^2+(Y_i-e_i)^2\bigr)
          +(A-1)Q,\qquad Q\ge0.
 \quad}
\end{equation}
Every solution has
\[
 (X_i,Y_i)=(\chi_A(e_i),\psi_A(e_i))
 \quad\text{for every }i.
\]
Conversely, these true principal pairs extend to a solution by fresh
anchor witnesses and an explicitly determined nonnegative integer $Q$.
\end{theorem}

\begin{proof}[Soundness]
Every term added to $W$ in \eqref{eq:one-row-oriented} is nonnegative,
so first $V\ge W$.  The anchor theorem is now applicable and gives
$V=\psi_A(W)$.  Only after this step classify each principal norm as
\[
 X_i=\sigma_i\chi_A(j_i),\qquad
 Y_i=\varepsilon_i\psi_A(j_i),\qquad j_i\ge0.
\]
The aggregate equation bounds $(Y_i-e_i)^2\le V-W<V$.  If
$j_i\ge B+1$, then, for either sign of $Y_i$,
\[
 |Y_i-e_i|\ge\psi_A(B+1)-e_i.
\]
Lemma~\ref{lem:next-index-square} contradicts the preceding strict
bound.  Consequently $0\le j_i\le B$ for every $i$.

Reduce the aggregate equation modulo $A-1$.  Since
$V=\psi_A(W)\equiv W$, \eqref{eq:one-row-residues} gives
\[
 A-1\mid E,
 \qquad
 E=\sum_{i=1}^m\bigl((\sigma_i-1)^2+
                         (\varepsilon_i j_i-e_i)^2\bigr).
\]
Here $(\sigma_i-1)^2$ is either $0$ or $4$.  Because $j_i\le B$ and
$1\le e_i\le B-1$, the absolute value of the other error is at most
$2B-1=W-1$.  Therefore
\begin{align*}
 0\le E&\le m\bigl(4+(W-1)^2\bigr)\\
       &\le(W-2)\bigl(4+(W-1)^2\bigr)<W^3<A-1.
\end{align*}
For completeness, the strict polynomial inequality used here is
\[
 W^3-(W-2)\bigl(4+(W-1)^2\bigr)=4W^2-9W+10>0.
\]
Thus $E=0$.  Each nonnegative summand vanishes, forcing
$\sigma_i=1$ and $\varepsilon_i j_i=e_i>0$.  This orients both
coordinates and identifies every index, including exclusion of $j_i=0$.
\end{proof}

\begin{proof}[Constructive completeness]
Take the true principal pairs and a complete anchor at $W$.  Write
$z=\psi_A(B-1)\ge1$.  Since $e_i\le B-1$ and $\psi_A(e_i)\ge e_i$,
the actual error sum satisfies
\begin{align*}
 E_{\rm actual}
 &=\sum_i\bigl((\chi_A(e_i)-1)^2+
                     (\psi_A(e_i)-e_i)^2\bigr)\\
 &\le m\bigl(\chi_A(B-1)^2+\psi_A(B-1)^2\bigr)
 \le m(A^2+1)z^2.
\end{align*}
On the other hand, $\psi_A(B)>Az$ and
$\chi_A(B)>(A-1)\psi_A(B)$, so
\[
 V=\psi_A(2B)=2\chi_A(B)\psi_A(B)>2(A-1)A^2z^2.
\]
It follows that
\[
 V-E_{\rm actual}
  >\bigl([2(A-1)-m]A^2-m\bigr)z^2>W.
\]
The last inequality uses $A>W+1$, which follows from $A-1>W^3$, and
$m\le W-2$.  Every $\chi_A(e_i)-1$ and $\psi_A(e_i)-e_i$ is divisible
by $A-1$, as is $V-W$.  Hence
\begin{equation}\label{eq:one-row-Q-witness}
 Q=\frac{V-W-E_{\rm actual}}{A-1}
\end{equation}
is a positive integer and supplies the aggregate equation.  The
inherited anchor completeness theorem supplies its private witnesses
afresh.  No pointwise transfer of unrelated old auxiliary tuples is
required.
\end{proof}

\subsection{A natural-domain specialization}

\begin{corollary}\label{cor:one-row-natural}
Under the same conditions on $A,W,m,e_i$, if all principal seed
coordinates are nonnegative, the simpler row
\begin{equation}\label{eq:one-row-natural}
 V=W+\sum_i(Y_i-e_i)^2+(A-1)Q,\qquad Q\ge0,
\end{equation}
already identifies the complete canonical pairs.
\end{corollary}
\begin{proof}
The preceding square bound gives $j_i\le B$.  Nonnegative principal
coordinates have positive abscissa and nonnegative ordinate signs.
The residue error is now $\sum_i(j_i-e_i)^2<W^3<A-1$, so the same
argument identifies the indices.  For completeness, put
$E_{\rm actual}=\sum_i(\psi_A(e_i)-e_i)^2\le m z^2$ and use
$V>2(A-1)z^2$.  Thus $V-W-E_{\rm actual}>0$, and division by $A-1$
gives the required natural $Q$.
\end{proof}

For a natural arithmetic source whose original common frame was
$W_0=2+\sum_{a\in\mathcal A}a$, choose the new frame $W=2W_0$.
Every positive seed exponent which occurs as an actual atom satisfies
$e_i<W/2$.  The change concerns the complete frame, including every
subsequent occurrence in its Pell certificate.  It is not a replacement
of one occurrence while leaving an old specialization elsewhere.

\subsection{The signed square frame and ordinary power outputs}

This subsection records an optional generalization.  The measured
378-integer equation instead uses the doubled linear frame of
Section~\ref{sec:frame}, with its atom signs proved first and a separate
lower parameter $L$.  The following square frame and signed-output
consequence are not included in its reported counts.

For a signed source one can use the complete common frame
\begin{equation}\label{eq:one-row-square-frame}
 W=2\left(2+\sum_{a\in\mathcal A}a^2\right).
\end{equation}
Its atoms are the actual semantic bases, exponents, and outputs before
their meanings have been established.  For every integer atom $k$,
\[
 W-4|k|\ge4+2k^2-4|k|=2(|k|-1)^2+2>0.
\]
Thus positive seed exponents $e_i=k_i$, and even selectively doubled
ones $e_i=2k_i$, lie below $W/2$.  Each claim must be traced to the
actual atom set; this inequality is not available for an arbitrary
unlisted expression.  The number of distinct positive seed atoms gives
the required count bound $m\le W-2$.

A polynomial large-frame certificate is
\[
 A^2-\bigl((W+1)^2-1\bigr)(WZ)^2=1,
 \qquad A\ge0,\quad Z\ge1.
\]
Pell classification at $W+1$ gives
$(A,WZ)=(\chi_{W+1}(h),\psi_{W+1}(h))$ for $h\ge1$.
The ordinate recurrence modulo $W$ gives $W\mid h$, so $h\ge W$.
Consequently
\[
 A\ge\chi_{W+1}(W)\ge(W+1)^W,
 \qquad A-1>W^3.
\]
This reasoning uses the literal frame and its norm, before interpreting
any power request.  It also supplies the stronger modulus bounds needed
for power recovery.

\begin{corollary}[Recovery of a signed power output]
\label{cor:one-row-signed-power}
Assume the strong frame bound $A\ge(W+1)^W$, and suppose a seed pair
has been oriented by Theorem~\ref{thm:one-row-oriented}.  Let
$1\le n<W$, $1\le e<W$, and $m\in\Z$ with $|m|<W$ be actual
base, exponent, and output atoms.  Put $T_n=2An-n^2-1$.  Then the
ordinary power congruence
\[
 \chi_A(e)+(n-A)\psi_A(e)=m+T_n h,\qquad h\in\Z,
\]
forces $m=n^e>0$.
\end{corollary}
\begin{proof}
The standard Pell recurrence modulo $T_n$ gives
$\chi_A(e)+(n-A)\psi_A(e)\equiv n^e\pmod{T_n}$: its characteristic
identity is $n^2-2An+1\equiv0$.  Therefore $T_n\mid(m-n^e)$.
For $1\le n<W<A$,
\[
 T_n\ge2A-2>W^W,
 \qquad |m-n^e|<W+W^{W-1}<W^W.
\]
The only possible multiple is zero.
\end{proof}

This corollary makes separate output squaring and sign selection
unnecessary at this interface.  It does not discard the source
conditions which establish that the actual bases and exponents are
positive.  Nor does it orient unrelated arithmetic-source coordinates.
An implementation must restore the complete original power-request
list and then check every later source use in its actual dependency
order.  For completeness the standard positive-power construction
provides the ordinary quotient; the corollary permits it to be signed.

\subsection{Exact row identities and essential hypotheses}

\begin{remark}[A useful replacement on the principal norm locus]
Let $\mathcal N_i=X_i^2-(A^2-1)Y_i^2-1$.  The free polynomial identity
\begin{align*}
 &(X_i-1)^2+(Y_i-e_i)^2\\
 &\quad-\bigl((AY_i)^2-2X_i+2-2e_iY_i+e_i^2\bigr)
 =\mathcal N_i
\end{align*}
allows the aggregate error to be replaced by the bracketed expression
when the principal norm rows are retained.  More precisely, the new
aggregate residual is the old one plus $\sum_i\mathcal N_i$, so the
complete integer zero set is unchanged.  If $AY_i$ already has an
actual stored multiplication wire, the replacement remains quadratic
in those coordinates.  Without that wire, writing $(AY_i)^2$ does
\emph{not} make it a quadratic polynomial in $A,Y_i$; a wire and its
defining row must be counted.  All size claims require the full emitted
graph, not this identity alone.
\end{remark}

The condition $Q\ge0$ is essential.  For an exact example, take
$W=4$, $A=66$, one requested index $e=1$, and the correct anchor at $4$.
Replace the seed by the positive Pell pair at $j=66=1+(A-1)$.
Both $\chi_A(66)-1$ and $\psi_A(66)-1$ are divisible by $65$,
and $V-4$ is divisible by $65$.  Thus
\[
 Q=\frac{V-4-(\chi_A(66)-1)^2-(\psi_A(66)-1)^2}{65}
\]
is an integer satisfying the aggregate equality.  It is negative,
because the wrong ordinate alone exceeds the proved cap.  Allowing an
unrestricted $Q$ would therefore admit a false index even though
$A-1=65>4^3$ and every other displayed condition holds.

Over unrestricted integers, the retained sign predicate $Q\ge0$ can
be realized by the familiar three-variable identity
\[
 Q=a^2+b^2+c(c+1),\qquad a,b,c\in\Z.
\]
If a quartic sum of quadratic residual squares is intended, keep $Q$
as an actual coordinate in $(A-1)Q$.  Directly substituting its
quadratic representation into that product would produce a cubic
residual.  Finally, the auxiliary letters in a straight-line evaluator
are abbreviations only when explicitly eliminated; their status in a
quadratic system must follow its stated variable list.

All arguments above use integer algebra, induction, Pell classification,
and finite divisibility.  They are formalizable within ZFC and assume
neither its consistency nor its inconsistency.  The new arithmetic
lemma supplies one part of the full syntactic-refutation equivalence;
it is not by itself a new end-to-end formal proof artifact.
